\documentclass{article}

\usepackage{amsfonts}
\usepackage{amsthm}
\usepackage{mathrsfs}
\usepackage{amsmath}
\usepackage{graphicx}
\usepackage{hyperref}
\usepackage{stackengine}
\usepackage[mathscr]{eucal}
\usepackage{listings}
\usepackage{algorithm}
\usepackage{algpseudocode}
\usepackage{float}
\usepackage{xcolor}
\usepackage{url}
\newtheorem{theorem}{Theorem}
\theoremstyle{definition}\newtheorem{definition}{DEF}
\theoremstyle{lemma}\newtheorem{lemma}{Lemma}
\theoremstyle{corrolary}\newtheorem{corrolary}{Cor}
\setlength{\parindent}{0pt}
\definecolor{codebg}{RGB}{248,248,248}
\definecolor{codekeyword}{RGB}{0,76,153}
\definecolor{codecomment}{RGB}{0,128,96}
\definecolor{codestring}{RGB}{163,21,21}
\definecolor{codenumber}{RGB}{102,102,102}
\lstdefinestyle{pythoncode}{
    language=Python,
    basicstyle=\ttfamily\small,
    keywordstyle=\color{codekeyword}\bfseries,
    commentstyle=\color{codecomment}\itshape,
    stringstyle=\color{codestring},
    numberstyle=\scriptsize\color{codenumber},
    numbers=left,
    numbersep=10pt,
    stepnumber=1,
    showstringspaces=false,
    breaklines=true,
    keepspaces=true,
    columns=fullflexible,
    tabsize=4,
    frame=single,
    framerule=0pt,
    rulecolor=\color{codebg},
    backgroundcolor=\color{codebg}
}
\lstset{style=pythoncode}

\begin{document}

\begin{lemma}
Let $X$ be a continuous random variable with CDF $F$.
Then $U= F(X)$ is a random variable with $U\sim Unif(0,1)$.
\end{lemma}
\begin{proof}
First, suppose $F$ is one-to-one and let $u\in [0,1]$.
Then there is an $x$ in the support $S_x$ such that $x=F^{-1}(u)$.
So the CDF is $F(F^{-1}(u))=u$.
Thus, $U\sim Unif(0,1)$.
In general, if $F$ is not one-to-one, we can define
\begin{equation*}
G(u)=\inf\{x : F(x)\ge u\}
\end{equation*}
Then $F(G(u))=u$.
So the distribution is $U\sim Unif(0,1)$.
\end{proof}
\begin{corrolary}
Let $X$ be a discrete random variable with CDF $F$ and support $S_x$.
Let $U\sim Unif(0,1)$ and $G(u)=\inf\{x : F(x)\ge u\}$.
Then $G(U)=X$.
\end{corrolary}

\begin{definition}
Monte Carlo is an algorithm that uses random number generation (RNG) to sample a random variable.
\end{definition}

\textbf{Note-}
Typically, we can not do true random number generation for a Monte Carlo algorithm.
Instead, we do pseudo-random number generation (pRNG) to do it and rely on the pRNG being "close to random."
This is because computers run deterministic programs, so generating truly random information is impossible in a computer.
pRNG algorithms usually require a seed, and it is common to use the time as a seed so repeated runs of an algorithm will not produce the same results.
In some specialized cases, an outside source of randomness is used to seed to have true RNG.
For example, Cloudflare maintains a wall of lava lamps in their office in San Francisco with a camera pointing at it.
Computers can then use this as a seed to generate truly random numbers.
However, using a good pRNG algorithm is much more common for everyday use.

\begin{definition}
Let $n, w, r\in \mathbb{N}$ such that $2^{nw-r}-1$ is a Mersenne prime.
Also let $1 \leq m < n$, let $0 \leq r < w$, let $1\leq u, s, t, l <w$, and let $a, b, c, f$ be binary sequences of length $w$.
A Mersenne-prime pRNG in the Mersenne Twister family can be defined by the following algorithm:
\begin{algorithm}[H]
\caption{Generic Mersenne-Prime Generation Step}
\begin{algorithmic}[1]
\State \textbf{Input:} Seed $x[0]$, state array $x[0 \dots n-1]$, index $k$
\State \textbf{Constants:} $n, w, m, r, a, b, c, f, u, s, t, l$
\State \textbf{Masks:} $\text{lower\_mask} \leftarrow 2^r - 1$, $\text{upper\_mask} \leftarrow (2^w - 1) - \text{lower\_mask}$
\If{the state is not initialized}
    \For{$i = 1$ to $n-1$}
        \State $x[i] \leftarrow \left(f \cdot \left(x[i-1] \text{ XOR } (x[i-1] \text{ RSHIFT } (w-2))\right) + i\right) \bmod 2^w$
    \EndFor
    \State $k \leftarrow n$
\EndIf
\If{$k \ge n$}
    \For{$i = 0$ to $n-1$}
        \State $y \leftarrow (x[i] \text{ AND } \text{upper\_mask}) \mid (x[(i+1) \bmod n] \text{ AND } \text{lower\_mask})$
        \State $xA \leftarrow y \text{ RSHIFT } 1$
        \If{$y \text{ AND } 1 \neq 0$}
            \State $xA \leftarrow xA \text{ XOR } a$
        \EndIf
        \State $x[i] \leftarrow x[(i + m) \bmod n] \text{ XOR } xA$
    \EndFor
    \State $k \leftarrow 0$
\EndIf
\State $y \leftarrow x[k]$
\State $k \leftarrow k + 1$
\State \Comment{Tempering}
\State $y \leftarrow y \text{ XOR } (y \text{ RSHIFT } u)$
\State $y \leftarrow y \text{ XOR } ((y \text{ LSHIFT } s) \text{ AND } b)$
\State $y \leftarrow y \text{ XOR } ((y \text{ LSHIFT } t) \text{ AND } c)$
\State $z \leftarrow y \text{ XOR } (y \text{ RSHIFT } l)$
\State \textbf{Output:} $z \in \{0, 1, \dots, 2^w-1\}$
\end{algorithmic}
\end{algorithm}
\end{definition}

\textbf{Note-}
The basic idea behind the Mersenne Twister is that it has 2 parts.
First, there is an internal state $x$ generated from a seed $x_0$.
This is generated by a recursive function called the Twister function.
At each step, the internal state is then transformed into the random number, called the tempering function.\\

\textbf{Ex-}
An example Mersenne Twister implementation in Python:
\lstinputlisting[firstline=20]{theory_of_probability/scripts/mersenne_twister.py}
An example of drawing numbers using this code with the MT19937 algorithm is provided in \path{theory_of_probability/scripts/run_mersenne_twister.py}

\begin{definition}
TestU01 is a library of statistical tests that determine whether a pRNG algorithm is distinguishable from draws of $Unif(0,1)$.
\end{definition}

\textbf{Ex-}
MT19937, the original implementation of the Mersenne Twister, succeeds on most tests in TestU01.
The 2 tests that it fails are related to binary representation of the generated numbers, which is due to the way in which it manipulates the binary data to generate the numbers.
Despite this, it is a common algorithm used in Monte Carlo methods for RNG.

\begin{theorem}
\textbf{Inverse Transform Method:}
Let $X$ be a random variable with CDF $F$ and let \verb|invF| be a function that computes $F^{-1}$.
Then the following generates samples from $X$ (in Python):
\end{theorem}
\begin{lstlisting}[language=Python]
import numpy as np

def invF(u: np.array) -> np.array:
    # Implementation of inverse CDF

u = np.random.rand(1)
x = invF(u)
\end{lstlisting}
\begin{proof}
Notice we use RNG to draw $U\sim Unif(0,1)$ as the variable \verb|u|.
Thus, by the previous lemma, \verb|x| is a sample from the desired distribution.
\end{proof}

\textbf{Ex-} Model rolling a 6-sided die:
\begin{lstlisting}[language=Python]
import numpy as np

def roll_die() -> np.array:
    u = np.random.rand(1)
    x = np.floor(6 * u) + 1
    return x
\end{lstlisting}

\begin{theorem}
Let $X$ be a random variable with CDF $X$ and finite variance, and $g: \mathbb{R}\rightarrow \mathbb{R}$ a function.
Then we can estimate the expected value, $\mathbb{E}(g(X))$, with:
\end{theorem}
\begin{lstlisting}[language=Python]
import numpy as np
from collections.abc import Callable

def expected_value(n: int, g: Callable[[np.array], np.array],
       invF: Callable[[np.array], np.array]) -> np.array:
    u = np.random.rand(n)
    x = invF(u)
    y = g(x)
    return np.mean(y)
\end{lstlisting}
\begin{proof}
First, notice we have sampled \verb|x| from $X\sim F(x)$ by the inverse transform method.
So \verb|y| is sampled from $g(X)$.
Let $\sigma^2=Var(g(X))$.
By Chebyshev's inequality,
$$\Pr\left( \middle| \frac{1}{n}\sum_k g(X_k) - \mathbb{E}(g(X))\middle| < \epsilon \right) \geq \frac{\sigma^2}{n^2}$$
Thus, for sufficiently large $n$, the average of \verb|y| is a good estimate of $\mathbb{E}(g(X))$.
\end{proof}

\textbf{Ex-}
Consider rolling 2 fair 20 sided dice and taking the maximum.
Estimate the average within 0.001 of the true value with 99\% probability.
Notice that the outcome of a fair, 20 sided die is between 1 and 20.
Thus, the outcome is a $\frac{19}{2}$-subgaussian random variable.
So we can use Hoeffding's inequality (Thm \ref{HoefdingInequality}) to bound the accuracy.
\begin{lstlisting}[language=Python]
import math
import numpy as np

epsilon = 0.001
delta = 0.01
n = math.ceil(19 ** 2 * math.log(2 / delta) / (2 * epsilon ** 2))
u = np.random.rand(n, 2)
dice = np.floor(20 * u) + 1
dice_max = np.amax(dice, axis=1)
print("Simulated rolling 1d20 with advantage, "
        + f"found mean of {np.mean(dice_max)} with margin of"
        + f"error {epsilon} and confidence {1-delta}.")
\end{lstlisting}

\begin{theorem}
\textbf{Box-Muller Transform:}
Let $U_1, U_2 {iid \atop \sim} Unif(0,1)$.
Let $Z_1 = \sqrt{-2\ln U_1}\cos (2\pi U_2)$ and $Z_2 = \sqrt{-2\ln U_1} \sin(2\pi U_2)$.
Then $Z_1, Z_2 \sim \mathcal{N}(0,1)$.
\end{theorem}
\begin{proof}
Let $R=\sqrt{-2\ln U_1}$.
Then $U_1 = e^{-R^2/2}$.
Thus, by the inverse transform method, the CDF of R is $F_R(r)=1-e^{-r^2/2}$.
So the PDF of $R$ is $f_R(r)=re^{-r^2/2}$.
Next, let $\phi = 2\pi U_2$.
So it has PDF $f(\phi)=\frac{1}{2\pi}$.
Clearly $R$ and $\phi$ are independent and have joint PDF $f(r,\phi)=\frac{r}{2\pi}e^{-r^2/2}$.
Also notice that $(R, \phi) \rightarrow (Z_1, Z_2)$ is a bijection with jacobian $R$.
Clearly, $R^2 = Z_1^2 + Z_2^2$.
Thus, the joint PDF of $Z$ is $f(z_1, z_2)=\frac{1}{2\pi} e^{-(z_1^2 + z_2^2)/2}$.
This is the PDF of $\mathcal{N}_2(0, \mathbb{I}_2)$.
Thus, we have the desired distribution.
\end{proof}

\begin{definition}
Rejection sampling is a technique in which a sample is taken from a larger space and rejected if it does not fall within the sample space.
\end{definition}

\textbf{Ex-}
In order to sample from the unit disk, we can sample a point $(x,y)\in [-1,1]^2$ and reject it if $x^2 + y^2 > 1$.\\

\textbf{Note-}
Rejection sampling may seem like a poor idea, especially in the case of sampling from a unit disk.
We could, instead, write an explicit formula to sample from the unit disk using the inverse transform method.
However, this transform would require the computation of multiple trigonometric functions, which can take a long time.
The probability of rejecting a point, on the other hand, is only $1-\frac{\pi}{4}$.
When timing the difference in low-level languages, such as C, rejection sampling is faster.
Python has unususally fast trigonometric functoins, so the reverse may be true for our example in \path{theory_of_probability/scripts/unit_disk_sampling.py}

\begin{definition}
Let $C\subseteq \mathbb{R}^n$.
Hit and run sampling is a sampling technique in which points are sampled recursively.
First, an initial point $x_0$ is the start of the algorithm.
Then, both a direction $d\in \mathbb{S}^{n-1}$ and the distance $A$ to the boundary of $C$ is computed.
Then, a distance $r\in [0,1]$ is sampled and $x_k=x_{k-1}+rd$.
\end{definition}

\textbf{Note-}
Hit and run sampling will sample uniformly from $C$ after an initial "warm up" period.
This is an example of the more general MCMC algorithm (\ref{MCMC}).

\end{document}
